\section{Постановка задачи}\label{ux433ux43bux430ux432ux430-1.-ux43fux43eux441ux442ux430ux43dux43eux432ux43aux430-ux437ux430ux434ux430ux447ux438}

\subsection{Актуальность и формулировка проблемы}\label{ux430ux43aux442ux443ux430ux43bux44cux43dux43eux441ux442ux44c-ux438-ux444ux43eux440ux43cux443ux43bux438ux440ux43eux432ux43aux430-ux43fux440ux43eux431ux43bux435ux43cux44b}

За последние пятнадцать лет онлайн-эксперименты стали одним из основных инструментов принятия решений в цифровых продуктах. В рекомендательных системах, рекламных платформах, маркетплейсах и системах ранжирования они применяются постоянно --- для выбора между альтернативными версиями интерфейсов, моделей и способов показа контента. Классический A/B-тест в этой среде остаётся базовым подходом к оценке изменений: при фиксированном случайном распределении трафика он даёт понятный и статистически интерпретируемый результат \cite{ref5}, \cite{ref6}, \cite{ref10}, \cite{ref11}. Вместе с тем по мере роста числа вариантов и ускорения продуктовых циклов проявляется недостаток фиксированного эксперимента --- пока накапливаются данные, часть пользователей продолжает видеть менее удачное решение. Поэтому в литературе всё чаще обсуждают алгоритмы многорукого бандита (MAB, multi-armed bandit) --- адаптивные методы, которые меняют распределение трафика уже в ходе эксперимента и стремятся быстрее направлять показы к более успешным вариантам \cite{ref1,ref2,ref3,ref4}, \cite{ref12}, \cite{ref13}.

В настоящей работе в качестве продуктовой задачи выбран рекомендательный сценарий --- выбор варианта показа пользователю; в качестве бизнес-метрики --- доля кликов (CTR, click-through rate). Такая постановка согласуется с открытым датасетом Open Bandit Dataset (OBD) и типичными задачами маркетплейсов, но не охватывает все продуктовые случаи: в других доменах целевой метрикой может быть конверсия, выручка, удержание и т. д.

A/B-тестирование и MAB решают родственные, но не тождественные задачи. A/B-тест нужен прежде всего для оценки эффекта и проверки статистической значимости различий. MAB нацелен на последовательный выбор действий так, чтобы одновременно продолжать обучение и получать как можно большую суммарную награду. В литературе это различие описывается через cumulative reward (накопленная награда) и regret (накопленное отставание от наилучшего действия) \cite{ref5}, \cite{ref12}, \cite{ref13}, \cite{ref18}, \cite{ref19}. Для рекомендательных систем с быстрой бинарной наградой цена неудачного показа выражается в потерянных кликах, конверсиях и доходе \cite{ref8}, \cite{ref12}, \cite{ref17}, \cite{ref18}.

Сводить вопрос к выбору «A/B или бандиты» как к единственной альтернативе некорректно: подходы решают разные задачи. Сопоставлять следует фиксированный эксперимент, ориентированный на надёжный статистический вывод, и адаптивный последовательный выбор, ориентированный на снижение потерь в процессе работы \cite{ref13}, \cite{ref18,ref19,ref20,ref21}. В научной литературе внимание уделяется не только алгоритмам, но и условиям применения, ограничениям статистического вывода, способам оценки политик на логах и переходу к контекстным бандитам, где выбор зависит от признаков пользователя.

Проблема, на которую направлена работа, имеет теоретический и прикладной уровни.

На теоретическом уровне A/B-тестирование и MAB оптимизируют разные цели: фиксированный эксперимент --- статистически корректный вывод об эффекте \cite{ref5}, \cite{ref19}, \cite{ref20}, MAB --- максимизацию награды и снижение regret в процессе обучения \cite{ref12}, \cite{ref13}. Работы Xiang и соавторов \cite{ref21} и Bleeker и Kaufmann \cite{ref22} отдельно подчёркивают, что выгода во время эксперимента и статистическая полезность итогового вывода --- не одно и то же. Для адаптивного назначения трафика стандартные p-value без специальных процедур некорректны \cite{ref19}, \cite{ref20}; для офлайн-сравнения политик требуется OPE (off-policy evaluation) с propensity \cite{ref6}, \cite{ref24}. Постановка «две оси --- оптимизация и статистический вывод» хорошо известна в литературе.

На прикладном уровне обзоры и интервью с практиками показывают, что эти различия не всегда переводятся в единую воспроизводимую схему принятия решений. Mattos и соавторы \cite{ref18} описывают организационные барьеры внедрения bandit-экспериментов; Dmitriev и соавторы \cite{ref11} --- типичные ошибки интерпретации метрик в контролируемых онлайн-экспериментах; Burtini и соавторы \cite{ref13} и Bleeker и Kaufmann \cite{ref22} дают критерии выбора по отдельности, но не объединяют их в единый проверяемый фреймворк. В реальной эксплуатации уже применяются и фиксированные A/B, и адаптивные схемы (в том числе как часть рекомендательного контура), однако правила интерпретации логов таких экспериментов часто смешивают «больше кликов во время теста» и «статистически значимый эффект». Разрыв между пользой онлайн-эксперимента с адаптивным назначением и корректностью статистического вывода по его логам --- центральная тема работы.

Цель курсовой --- систематизировать известный теоретический и прикладной разрыв в воспроизводимом коде: сначала эмпирически показать, когда адаптивное назначение трафика приносит измеримую пользу в онлайн-эксперименте, затем отдельно проверить, какие процедуры дают статистически обоснованный вывод и где стандартные приёмы дают сбой. 

\subsection{Бизнес-контекст и продуктовое решение}\label{ux431ux438ux437ux43dux435ux441-ux43aux43eux43dux442ux435ux43aux441ux442-ux438-ux43fux440ux43eux434ux443ux43aux442ux43eux432ux43eux435-ux440ux435ux448ux435ux43dux438ux435}

Рассматривается типовой сценарий цифрового продукта с рекомендательным или ранжирующим интерфейсом: e-commerce, маркетплейс или рекламная платформа. На каждом показе система выбирает один из нескольких вариантов --- товар, креатив или позицию в выдаче --- и наблюдает бинарный отклик пользователя (клик или его отсутствие). Продуктовая команда регулярно запускает эксперименты, чтобы выбрать лучший вариант и одновременно не терять трафик на заведомо слабые альтернативы.

Центральное продуктовое решение, на которое должна помочь ответить курсовая, формулируется так: при оптимизации CTR в задаче с быстрой бинарной наградой и несколькими альтернативами --- когда целесообразно использовать адаптивное MAB-распределение трафика вместо фиксированного разбиения трафика, и когда A/B-тестирование остаётся предпочтительным для статистически надёжного вывода?

В рамках работы сравниваются четыре способа назначения трафика. В качестве бейзлайна выступает \texttt{fixed\_ab} --- равномерное или заранее зафиксированное распределение между вариантами, соответствующее базовой схеме A/B-теста. Три MAB-альтернативы: $\varepsilon$-greedy (с вероятностью $\varepsilon$ --- случайный вариант, иначе --- лучший по эмпирике), UCB1 (выбор по upper confidence bound) и Thompson Sampling (байесовский сэмплинг из апостериорных распределений). Продуктовое решение принимается по двум независимым критериям: по метрикам последовательной оптимизации (reward, regret) и по корректности статистического вывода (ошибка I рода, мощность, p-value). Смешивать эти критерии в одном заключении нельзя \cite{ref13}, \cite{ref21}.

\textbf{Обозначения} (имена политик совпадают с кодом в \texttt{src/bandits/}):

\begin{itemize}
\tightlist
\item
  \texttt{fixed\_ab} --- фиксированное разбиение трафика (FixedABPolicy): вероятности показа вариантов не меняются по ходу эксперимента; при двух руках --- разбиение 50/50;
\item
  \texttt{epsilon\_greedy} ($\varepsilon$-greedy) --- с вероятностью $\varepsilon$ выбирает случайный вариант, иначе --- лучший по текущей статистике;
\item
  \texttt{ucb1} --- Upper Confidence Bound;
\item
  \texttt{thompson\_sampling} (TS) --- сэмплирование Томпсона;
\item
  \texttt{linucb} --- контекстный бандит LinUCB.
\end{itemize}

\textbf{Прочие обозначения:} T --- горизонт (число шагов/событий); seed --- начальное значение генератора псевдослучайных чисел; regret --- накопленные недобранные клики относительно оракула; suboptimal\_share --- доля шагов с показом не лучшего варианта; propensity --- вероятность назначения варианта при политике сбора логов (см. сокращения); доля совпадений с логом --- при OPE доля событий, где политика-кандидат выбирает ту же руку, что записана в логе (см. сокращения); IPS --- inverse propensity scoring; SNIPS --- self-normalized IPS; ESS (effective sample size) --- эффективный размер выборки весов; OPE (off-policy evaluation) --- офлайн-оценка политики-кандидата по логам; BTS --- Bernoulli Thompson Sampling, адаптивная политика сбора в OBD; CTR\_A, CTR\_B --- CTR контрольного и тестового вариантов.

\subsection{Участники и сценарии использования}\label{ux443ux447ux430ux441ux442ux43dux438ux43aux438-ux438-ux441ux446ux435ux43dux430ux440ux438ux438-ux438ux441ux43fux43eux43bux44cux437ux43eux432ux430ux43dux438ux44f}

Результаты работы ориентированы на продуктового менеджера, выбирающего между A/B и MAB для следующего эксперимента по CTR; ML- и RecSys-инженера, внедряющего политику назначения трафика и офлайн-оценку на логах; аналитика экспериментов, интерпретирующего p-value и доверительные интервалы без смешивания метрик оптимизации и статистического вывода; научного руководителя и рецензента, проверяющих методологическую корректность и воспроизводимость пайплайна.

Типичные сценарии: запуск онлайн-эксперимента с несколькими вариантами показа, товара или креатива; офлайн-оценка гипотетической политики на исторических логах с известными propensity; подготовка обоснования выбора между фиксированным разбиением трафика и адаптивным назначением трафика для продуктовой презентации или внутреннего обзора.

\subsection{Цена ошибки}\label{ux446ux435ux43dux430-ux43eux448ux438ux431ux43aux438}

В задаче оптимизации CTR цена ошибки выражается в потерянных кликах, конверсиях и выручке, а не в риске для жизни. Это допускает компромисс между качеством оптимизации, простотой интерпретации и стоимостью внедрения, но не снимает требований к статистической корректности вывода.

Ошибка чрезмерной эксплуатации (false exploitation) --- система слишком рано закрепляется на слабом варианте и продолжает показывать его пользователям, хотя существует лучшая альтернатива. Ошибка чрезмерного исследования (false exploration) --- из-за высокого $\varepsilon$ или неудачной политики исследования слишком долго тратится трафик на заведомо слабые варианты. Обе ведут к потерям в кликах и GMV; первая аргументирует применение MAB для снижения regret, вторая --- необходимость абляции гиперпараметров.

На уровне статистического вывода отдельно рассматриваются ложноположительные и ложноотрицательные решения. Риск ложноположительного вывода при адаптивном назначении трафика следует из литературы по последовательному анализу и always-valid inference (см. сокращения) \cite{ref19}, \cite{ref20}: стандартный двухвыборочный тест предполагает фиксированный дизайн и независимость наблюдений, которые нарушаются при MAB-назначении трафика. Xiang и соавторы \cite{ref21} рассматривают MAB vs A/B с точки зрения доверительных интервалов и мощности, подтверждая, что статистический вывод на adaptive-логах --- отдельная задача. Ложноотрицательный вывод --- когда реальный, хотя и небольшой эффект не обнаруживается из-за короткого горизонта или малого MDE --- обосновывает планирование sample size \cite{ref5} и не должно смешиваться с метриками regret.

\textbf{Таблица 1.1 --- Типы ошибок в постановке CTR}

\begingroup\footnotesize\setlength{\tabcolsep}{4pt}
\begin{xltabular}{\linewidth}{@{}
  >{\raggedright\arraybackslash}Y
  >{\raggedright\arraybackslash}Y
  >{\raggedright\arraybackslash}Y
  >{\raggedright\arraybackslash}Y@{}}
\toprule\noalign{}
\midrule\noalign{}
\endhead
\bottomrule\noalign{}
\endlastfoot
Чрезмерная эксплуатация & Показы слабого варианта при наличии лучшего & Потерянные клики, GMV & Аргумент за MAB: снижение regret \\
Чрезмерное исследование & Высокий $\varepsilon$, долгое исследование слабых вариантов & Аналогичные потери & Абляция гиперпараметров ($\varepsilon$) \\
Ложноположительный вывод & Наивный A/B на adaptive-логах & Ложное принятие гипотезы & Отдельный слой статистического вывода; фиксированное разбиение трафика \\
Ложноотрицательный вывод & Малый MDE при коротком горизонте & Упущенное улучшение & Планирование sample size; не путать с regret \\
\end{xltabular}
\endgroup

\subsection{Ограничения, альтернативы и область применения}\label{ux43eux433ux440ux430ux43dux438ux447ux435ux43dux438ux44f-ux430ux43bux44cux442ux435ux440ux43dux430ux442ux438ux432ux44b-ux438-ux43eux431ux43bux430ux441ux442ux44c-ux43fux440ux438ux43cux435ux43dux435ux43dux438ux44f}

Работа ограничена по масштабу. Не рассматриваются нейросетевые бандиты и работа рекомендательных моделей в продакшене. \textbf{Batch-режим} --- пакетная обработка логов: политика обновляется не после каждого события, а после накопления пакета (\texttt{batch\_size} > 1); в курсовой реализован как упрощённый симулятор на исторических логах и не заменяет строгую контрфактическую оценку (replay-OPE с IPS/SNIPS). Контекстные бандиты представлены теоретически и в виде демонстрации на игрушечных данных (LinUCB), но не как основной эмпирический результат. Для реальных данных используется учебная подвыборка OBD (10 000 записей) наряду с полным random/all в отдельных экспериментах (E9a).

Среди альтернатив: классический A/B с фиксированным горизонтом и корректным статистическим выводом \cite{ref5}, \cite{ref19}, \cite{ref20}; CUPED для снижения дисперсии fixed A/B \cite{ref7} --- в курсовой проверен на синтетике (E15); OPE на логах (IPS, SNIPS, DR) \cite{ref6}, \cite{ref24}, \cite{ref27}; contextual bandit pipeline \cite{ref15}, \cite{ref16}.

Область применимости --- задачи с быстрой бинарной наградой (CTR), умеренным числом рук (от двух до нескольких десятков), наличием логов с propensity для офлайн-оценки и продуктовой культурой, в которой команда различает оптимизацию и статистический вывод.

\subsection{Исследовательские вопросы и критерии успеха}\label{ux438ux441ux441ux43bux435ux434ux43eux432ux430ux442ux435ux43bux44cux441ux43aux438ux435-ux432ux43eux43fux440ux43eux441ux44b-ux438-ux43aux440ux438ux442ux435ux440ux438ux438-ux443ux441ux43fux435ux445ux430}

Исследование строится вокруг восьми направлений (RQ1--RQ7 и подвопросы RQ3.1--RQ3.3 по статистическому выводу). 

\textbf{RQ1:} снижают ли MAB-политики cumulative regret относительно \texttt{fixed\_ab}? 

\textbf{RQ2:} как различаются политики при OPE на реальных логах? 

\textbf{RQ3.1} (адаптивные логи): сохраняется ли контроль ошибки I рода при наивном A/B? 

\textbf{RQ3.2:} восстанавливает ли IPS-weighted inference корректность? 

\textbf{RQ3.3:} подглядывание за промежуточными результатами при фиксированном разбиении трафика и последовательные процедуры (OBF, mSPRT) \cite{ref19,ref33}? 

\textbf{RQ4:} компромисс $\varepsilon$. 

\textbf{RQ5:} контекстный бандит vs неконтекстный. 

\textbf{RQ6:} устойчивость OPE к политике логирования и сегменту. 

\textbf{RQ7:} классическая A/B-постановка (два варианта, гипотеза CTR\_B \textgreater{} CTR\_A).

Соответствие вопросов и экспериментов --- в табл. 1.2; сводка параметров прогонов --- в табл. 3.1 (§3.5).

\textbf{Таблица 1.2 --- Исследовательские вопросы и эксперименты}

\begingroup\footnotesize\setlength{\tabcolsep}{4pt}
\begin{xltabular}{\linewidth}{@{}
  >{\raggedright\arraybackslash}Y
  >{\raggedright\arraybackslash}Y
  >{\raggedright\arraybackslash}Y@{}}
\toprule\noalign{}
\midrule\noalign{}
\endhead
\bottomrule\noalign{}
\endlastfoot
RQ1 & Regret MAB vs \texttt{fixed\_ab} & E1 \\
RQ2 & OPE на логах с propensity & E2, E2b/c, E9a, E9b \\
RQ3.1 & Наивный статистический вывод на adaptive-логах & E4, E5 \\
RQ3.2 & IPS-weighted inference & E11 \\
RQ3.3 & Подглядывание / последовательные (OBF, mSPRT) & E14 \\
RQ4 & Абляция $\varepsilon$ (exploration--exploitation) & E6 \\
RQ5 & Контекстный бандит vs неконтекстный & E8b \\
RQ6 & Устойчивость OPE к политике логирования и сегменту & E2b/c, E9b \\
RQ7 & Классическая A/B-постановка (два варианта) & E12, E13 \\
\end{xltabular}
\endgroup

Критерии успешного выполнения: воспроизводимый репозиторий с тестами, README и зафиксированными seeds; бейзлайн (\texttt{fixed\_ab}) и не менее трёх MAB-альтернатив; двухуровневая эмпирика на синтетике и OBD; раздельные выводы по метрикам оптимизации и статистического вывода; сформулированные ограничения и продуктовая рекомендация. Количественные результаты приводятся в главе 3 (§3.6--3.7).

\subsection{Цель, задачи и границы работы}\label{ux446ux435ux43bux44c-ux437ux430ux434ux430ux447ux438-ux438-ux433ux440ux430ux43dux438ux446ux44b-ux440ux430ux431ux43eux442ux44b}

\textbf{Цель} --- на основе обзора литературы \cite{ref12}, \cite{ref13}, \cite{ref18}, \cite{ref21}, \cite{ref22} подготовить воспроизводимый фреймворк сравнения \texttt{fixed\_ab} и MAB-политик в задачах CTR и рекомендаций, в которой разведены последовательная оптимизация и статистический вывод.

\textbf{Задачи работы:}

\begin{enumerate}
\def\labelenumi{\arabic{enumi}.}
\tightlist
\item
  Реализовать бейзлайн (\texttt{fixed\_ab}), MAB-алгоритмы ($\varepsilon$-greedy, UCB1, Thompson Sampling) и LinUCB в едином воспроизводимом пайплайне.
\item
  Сформулировать и проверить RQ1--RQ7 и RQ3.1--RQ3.3 (§1.6) на синтетических данных, контекстной синтетике и OBD.
\item
  Разделить в программной реализации политику назначения трафика и процедуру A/B inference.
\item
  Провести офлайн-оценку политик на логах с propensity (IPS, SNIPS) и отдельное исследование корректности наивного статистического вывода при адаптивном назначении (E11).
\item
  Сформулировать продуктовую рекомендацию и ограничения с опорой на литературу и эмпирические результаты главы 3 (§3.6--3.7).
\end{enumerate}

\textbf{Границы работы:} новые bandit-алгоритмы или промышленная система реального времени; первое сравнение политик логирования на ZOZOTOWN (уже выполнено в OBD \cite{ref26}); универсальное доказательство превосходства MAB над A/B \cite{ref18}, \cite{ref21}, \cite{ref22}; CUPED \cite{ref7} --- E15 на синтетике; DR \cite{ref27} --- в обзоре и §4.3 без реализации.

\textbf{Научная новизна и вклад.} Теоретическое различие оптимизации и статистического вывода, риски наивного вывода на adaptive-логах известны \cite{ref5,ref13,ref19,ref21}; OBD \cite{ref26} даёт воспроизводимый OPE-бенчмарк, но не связывает в одном фреймворке пользу MAB и карту процедур статистического вывода. Вклад курсовой --- воспроизводимый двухблочный фреймворк (блок I: regret/клики; блок II: ошибка I рода, мощность, IPS, последовательные процедуры) с эмпирической калибровкой на синтетике и OBD и продуктовой матрицей решений (гл. 4, табл. 4.1). Новые алгоритмы не предложены; работа систематизирует и проверяет известные тезисы в постановке CTR с открытым кодом. Количественные результаты --- глава 3; итоги и возможные направления развития --- глава 4.

\subsection{Четыре типа оценки}\label{ux447ux435ux442ux44bux440ux435-ux442ux438ux43fux430-ux43eux446ux435ux43dux43aux438}

В работе используются четыре независимых типа оценки \cite{ref12}, \cite{ref13}, \cite{ref21}; результаты одного не подменяют результаты другого. \textbf{Онлайн-оптимизация} --- сколько кликов накоплено за горизонт (reward, regret, suboptimal\_share); эксперименты E1, E6, E8b, E12. \textbf{Статистический вывод} --- корректность проверки гипотезы о превосходстве варианта B (p-value, rejection rate, мощность); E4, E5, E11, E14, блок inference в E12. \textbf{Офлайн-оценка (OPE)} --- оценка политики-кандидата по логам с propensity (SNIPS, ESS, доля совпадений с логом); E2, E2b/c, E9, E13 (иллюстрация pairwise). Разведочные batch-прогоны E3 и E8 (пакетная обработка логов, см. §1.5) --- приложение А, не IPS/SNIPS replay-OPE. \textbf{Продуктовая применимость} --- время задержки одного решения алгоритма; P1.

Подробное описание метрик и группировка результатов по типам оценки --- в §3.2--3.6.

\section{Обзор литературы}\label{ux433ux43bux430ux432ux430-2.-ux43eux431ux437ux43eux440-ux43bux438ux442ux435ux440ux430ux442ux443ux440ux44b}

\subsection{Введение к обзору}\label{ux432ux432ux435ux434ux435ux43dux438ux435-ux43a-ux43eux431ux437ux43eux440ux443}

Обзор служит теоретической основой для экспериментальной части. Центральная ось --- онлайн-эксперимент в рекомендательных и ранжирующих системах: как в индустрии устроены известные продовые решения, как из них получают статистически значимый результат и где возникают слабые места при интерпретации логов. Сначала рассматриваются классические A/B-тесты и промышленные практики контролируемых экспериментов, затем --- адаптивные схемы (MAB) как часть рекомендательного контура. После этого анализируются типичные ошибки смешивания оптимизации и статистического вывода, контекстные бандиты и эмпирические основания (OBD). В конце формулируется постановка курсовой: MAB --- не полная замена A/B, а инструмент онлайн-эксперимента с адаптивным назначением; статистический вывод по таким логам требует отдельных процедур.

\subsection{Классическое A/B-тестирование как основа онлайн-экспериментов}\label{ux43aux43bux430ux441ux441ux438ux447ux435ux441ux43aux43eux435-ab-ux442ux435ux441ux442ux438ux440ux43eux432ux430ux43dux438ux435-ux43aux430ux43a-ux43eux441ux43dux43eux432ux430-ux43eux43dux43bux430ux439ux43d-ux44dux43aux441ux43fux435ux440ux438ux43cux435ux43dux442ux43eux432}

Под онлайн-экспериментом обычно понимается контролируемое изменение части цифровой системы, при котором пользователи случайным образом распределяются между несколькими вариантами, а затем по заранее выбранным метрикам оценивается влияние вмешательства на поведение аудитории \cite{ref5}, \cite{ref10}. Эксперимент --- способ получить обоснованный вывод о последствиях изменения в реальной среде, а не только на офлайн-моделях или экспертных предположениях \cite{ref5}, \cite{ref20}.

Классическое A/B-тестирование остаётся наиболее распространённой формой такого эксперимента. В простейшем случае пользователи делятся на контрольную группу, которая видит текущий вариант системы, и тестовую --- с изменённым вариантом. После этого сравниваются значения одной или нескольких метрик: CTR, конверсии, удержания, выручки \cite{ref5}, \cite{ref10}, \cite{ref11}. Теоретическая сила A/B-теста в том, что при случайном и фиксированном распределении трафика различия между группами можно интерпретировать как следствие экспериментального воздействия при соблюдении базовых предпосылок статистического анализа \cite{ref5}.

В работах Kohavi и соавторов A/B-тест описывается как целостная исследовательская процедура: постановка гипотезы, выбор основной метрики, определение длительности эксперимента, контроль качества данных и корректная интерпретация результата \cite{ref5}, \cite{ref10}. Ошибки возникают не только на уровне формулы, но и на уровне постановки: слабая связь метрики с целью продукта, краткосрочный рост локального показателя при ухудшении системы в целом, сезонность, сбои логирования, пересечение экспериментов \cite{ref6}, \cite{ref10}, \cite{ref11}. Статистически значимый результат может оказаться практически бессмысленным с точки зрения продукта \cite{ref5}, \cite{ref10}, \cite{ref11}.

Отдельное направление связано с повышением чувствительности экспериментов. В работе Deng и соавторов по CUPED показано, что использование доэкспериментальных данных уменьшает разброс оценок и повышает мощность теста без дополнительного трафика \cite{ref7}. Даже в рамках фиксированного эксперимента существует развитый набор приёмов: уменьшение дисперсии, разделение целевых и защитных метрик, планирование с учётом мощности и MDE \cite{ref5}, \cite{ref7}, \cite{ref10}, \cite{ref20}.

Существенное значение имеет способ наблюдения за результатами во времени. Обычные p-value нельзя без оговорок использовать при произвольной остановке эксперимента; постоянное обращение к промежуточным результатам увеличивает риск ложноположительных выводов \cite{ref19}, \cite{ref20}. Под фиксированным горизонтом понимается ситуация, когда объём выборки или длительность заданы заранее и решение принимается после достижения установленного объёма данных. Johari и соавторы \cite{ref19} показывают, что для непрерывного мониторинга нужны специальные методы always-valid inference (корректный статистический вывод при непрерывном мониторинге без завышения ошибки I рода; см. сокращения). A/B-тестирование следует рассматривать как метод построения статистически надёжного вывода о различии между вариантами, а не только как способ случайного распределения пользователей.

\subsubsection{Продовые решения в рекомендациях и онлайн-экспериментах}\label{ux43fux440ux43eux434ux43eux432ux44bux435-ux440ux435ux448ux435ux43dux438ux44f-ux432-ux440ux435ux43aux43eux43cux435ux43dux434ux430ux446ux438ux44fux445-ux438-ux43eux43dux43bux430ux439ux43d-ux44dux43aux441ux43fux435ux440ux438ux43cux435ux43dux442ux430ux445}

В индустрии сложились несколько устойчивых схем, которые в той или иной форме решают задачу «какой вариант показывать пользователю» и «как доказать, что изменение работает». Их важно различать: одни относятся к назначению трафика во время эксперимента, другие --- к выводу о статистической значимости после его завершения (или по ходу при специальных процедурах).

\textbf{Фиксированный A/B как эталон проверки гипотезы.} На зрелых экспериментальных платформах (описание процесса --- у Kohavi и соавторов \cite{ref5}, \cite{ref10}, \cite{ref20}) стандартный цикл выглядит так: формулируется гипотеза и основная метрика (CTR, конверсия, выручка); заранее планируется длительность или объём выборки с учётом MDE и мощности; пользователи случайно и независимо от промежуточных результатов распределяются между контролем и тестом; по достижении заранее заданного горизонта применяется двухвыборочный тест или эквивалентная процедура. Статистически значимый результат здесь означает: при истинном отсутствии эффекта доля ложных срабатываний контролируется уровнем $\alpha$ (обычно 5\%), а при наличии эффекта заданного размера тест обладает достаточной мощностью. Дополнительные приёмы --- снижение дисперсии (CUPED \cite{ref7}), разделение целевых и защитных метрик, аудит качества логов \cite{ref6}, \cite{ref10}, \cite{ref11}.

\textbf{Адаптивное назначение (MAB) как часть рекомендательного контура.} В задачах с быстрой бинарной наградой и множеством альтернатив (товары, креативы, позиции) фиксированное разбиение трафика тратит заметную долю показов на слабые варианты. Scott \cite{ref12} и обзор Burtini et al.~\cite{ref13} описывают MAB-эксперименты в онлайн-сервисах как способ снизить regret во время обучения: $\varepsilon$-greedy, UCB и Thompson Sampling перераспределяют трафик к вариантам с лучшим промежуточным CTR. В рекомендациях контекстные бандиты (LinUCB и родственные методы) использовались для персонализации новостной ленты Yahoo \cite{ref6}; на уровне продуктовых кейсов аналогичная логика --- быстрая адаптация к отклику пользователя вместо цикла «обучили офлайн --- зафиксировали разбиение --- дождались конца теста» \cite{ref23}. В таких схемах польза измеряется накопленной наградой (клики, конверсии), а не p-value по сравнению двух групп.

\textbf{Офлайн-оценка политик по логам онлайн-эксперимента.} Когда эксперимент уже проведён под одной политикой логирования, сравнить гипотетических кандидатов без нового трафика можно через off-policy evaluation (IPS, SNIPS, DR) при наличии propensity \cite{ref6}, \cite{ref24}, \cite{ref27}, \cite{ref30}. Open Bandit Dataset \cite{ref26} --- воспроизводимый пример такого пайплайна на fashion e-commerce: логи с random и адаптивным (BTS) назначением, многорукая постановка, готовые propensity. Для debiasing в рекомендациях важна постановка «рекомендация как treatment» \cite{ref31}. Это не заменяет A/B для проверки гипотезы, но позволяет ранжировать политики на исторических данных при корректной постановке.

\textbf{Типичная продуктовая комбинация.} На практике редко выбирают «только бандит» или «только A/B». Bleeker и Kaufmann \cite{ref22} и Xiang et al.~\cite{ref21} подчёркивают: MAB снижает стоимость эксперимента в смысле накопленной награды, тогда как для принятия гипотезы о превосходстве B по-прежнему нужен защищённый статистический вывод --- отдельный fixed A/B, sequential design при мониторинге \cite{ref19} или специализированный анализ adaptive-логов. Mattos et al.~\cite{ref18} добавляют организационный слой: даже при известных алгоритмах команды путают метрики оптимизации и статистического вывода, что усиливает риск ложного принятия гипотезы.

\textbf{Таблица 2.2 --- Продовые схемы: цель, как получают «значимость», слабое место}

\begingroup\footnotesize\setlength{\tabcolsep}{4pt}
\begin{xltabular}{\linewidth}{@{}
  >{\raggedright\arraybackslash}Y
  >{\raggedright\arraybackslash}Y
  >{\raggedright\arraybackslash}Y
  >{\raggedright\arraybackslash}Y@{}}
\toprule\noalign{}
\midrule\noalign{}
\endhead
\bottomrule\noalign{}
\endlastfoot
Fixed A/B & Проверка гипотезы, guardrail-метрики & p-value / ДИ при фиксированном горизонте, план мощности \cite{ref5} & Подглядывание без поправки \cite{ref19}; слабая метрика \cite{ref11} \\
MAB (TS, $\varepsilon$-greedy, UCB) & Больше кликов во время теста & Regret, cumulative reward \cite{ref12}, \cite{ref13} & Наивный A/B по логам бандита → завышенная ошибка I рода \cite{ref19}, \cite{ref21} \\
Contextual bandit & Персонализация показа \cite{ref6}, \cite{ref23} & Онлайн-regret; офлайн --- OPE с propensity & Высокая стоимость работы в продакшене; batch-OPE без строгой propensity \\
OPE на логах & Сравнение кандидатов без нового трафика \cite{ref24}, \cite{ref26} & SNIPS, ESS, доля совпадений с логом & Низкий ESS на adaptive-логах; шумное ранжирование \\
Sequential / always-valid & Мониторинг фиксированного A/B по ходу \cite{ref19} & OBF, mSPRT с контролем ошибки I рода & Не исправляет адаптивное назначение бандита \\
\end{xltabular}
\endgroup

\subsubsection{Слабые места интерпретации логов онлайн-эксперимента}\label{ux441ux43bux430ux431ux44bux435-ux43cux435ux441ux442ux430-ux438ux43dux442ux435ux440ux43fux440ux435ux442ux430ux446ux438ux438-ux43bux43eux433ux43eux432-ux43eux43dux43bux430ux439ux43d-ux44dux43aux441ux43fux435ux440ux438ux43cux435ux43dux442ux430}

Литература выделяет несколько повторяющихся ошибок, когда логи онлайн-эксперимента читают «как в учебнике по A/B», хотя дизайн этому не соответствует.

\textbf{Смешивание оптимизации и статистического вывода.} Высокий cumulative CTR или низкий regret у MAB не доказывают статистическую значимость превосходства B над A: трафик перераспределён в пользу лучшего варианта, выборки в группах несимметричны, наблюдения зависимы \cite{ref13}, \cite{ref21}. Dmitriev et al.~\cite{ref11} перечисляют родственные ловушки на уровне метрик: локальный рост CTR при ухудшении системы в целом, неверный выбор primary metric.

\textbf{Наивный статистический вывод на adaptive-логах.} Стандартный z-критерий предполагает фиксированное разбиение трафика и независимость условных на успех наблюдений. Логи Thompson Sampling или $\varepsilon$-greedy этим предпосылкам не удовлетворяют: rejection rate при null завышается, мощность искажается \cite{ref19}, \cite{ref20}, \cite{ref21}. Для постфактумного вывода на таких логах нужны IPS-weighted процедуры с propensity \cite{ref6}, \cite{ref24} --- с оговоркой о мощности при редких показах слабой руки.

\textbf{Подглядывание за промежуточными результатами при фиксированном разбиении трафика.} Даже при честном 50/50 многократный просмотр промежуточных p-value и остановка при первом «значимом» результате раздувает ошибку I рода \cite{ref19}, \cite{ref20}. Индустриальные платформы либо фиксируют горизонт заранее \cite{ref5}, либо применяют последовательный групповой дизайн (O'Brien--Fleming) либо тесты с корректным выводом при непрерывном мониторинге (always-valid, см. сокращения) \cite{ref19}. Это другая проблема, чем адаптивное назначение бандита: последовательные процедуры не заменяют IPS на bandit-логах.

\textbf{OPE без учёта политики логирования.} Оценка политики-кандидата на логах, собранных адаптивной политикой (BTS в OBD), даёт низкий ESS и нестабильное ранжирование \cite{ref26}. Вывод «лучший алгоритм по SNIPS» на малой выборке при низком CTR может быть артефактом дисперсии, а не продуктовым открытием \cite{ref24}.

Эти слабые места --- неверная интерпретация логов конкретного онлайн-эксперимента --- составляют мотивацию экспериментальной части курсовой (глава 3).

\subsection{Компромисс между A/B и MAB в литературе}\label{ux43aux43eux43cux43fux440ux43eux43cux438ux441ux441-ux43cux435ux436ux434ux443-ab-ux438-mab-ux432-ux43bux438ux442ux435ux440ux430ux442ux443ux440ux435}

Сопоставление подходов показывает, что A/B-тестирование и MAB оптимизируют разные цели. Fixed A/B ориентирован на статистически корректный вывод об эффекте \cite{ref5}, \cite{ref19}, \cite{ref20}; MAB --- на максимизацию награды и снижение regret \cite{ref12}, \cite{ref13}. Xiang и соавторы \cite{ref21} формулируют это как различие между выгодой во время эксперимента и полезностью итогового статистического заключения. Bleeker и Kaufmann \cite{ref22} рассматривают MAB как способ снизить стоимость эксперимента, но не как полную замену A/B. В реальной эксплуатации нередко нужны оба типа результата --- больше кликов при обучении и надёжное принятие гипотезы о превосходстве B --- однако для этого требуются разные метрики и разные процедуры анализа \cite{ref13}, \cite{ref21}, а не один «универсальный» показатель вроде cumulative CTR или единого p-value.

С точки зрения стоимости внедрения классический A/B остаётся относительно дешёвым там, где уже есть зрелая экспериментальная платформа \cite{ref5}, \cite{ref20}. MAB требует адаптивного назначения, мониторинга и аккуратной организации логирования; Mattos и соавторы \cite{ref18}, опираясь на интервью с экспериментаторами из пяти компаний, показывают, что организационная цена может превышать чисто алгоритмическую. По времени задержки простые MAB-политики ($\varepsilon$-greedy, UCB1, TS) не уступают фиксированному разбиению; тяжёлые контекстные модели существенно дороже \cite{ref15}, \cite{ref16}.

Интерпретируемость --- один из главных барьеров MAB \cite{ref18}: проще объяснить «группа A vs B, p-value, доверительный интервал», чем динамику regret и смещение трафика во времени. Офлайн-оценка для A/B в простейшем виде сводится к сравнению групп на логах фиксированного эксперимента; для контрфактического сравнения гипотетических MAB-политик требуется OPE с propensity \cite{ref6}, \cite{ref24}.

При росте числа рук фиксированное разбиение становится дорогим по трафику: для статистически устойчивого вывода на каждый вариант нужен больший объём данных \cite{ref5}. MAB лучше масштабируется в многоруких сценариях, перераспределяя показы к более удачным альтернативам \cite{ref13} --- что заметно при десятках товаров или креативов в e-commerce \cite{ref21}, \cite{ref26}.

Литература \cite{ref11}, \cite{ref12}, \cite{ref19}, \cite{ref21}, \cite{ref22} описывает типичные ошибки смешивания целей: считать cumulative CTR MAB доказательством превосходства B над A (оптимизация $\neq$ статистический вывод); применять z-test к логам adaptive-политики без коррекции; сравнивать MAB и A/B одним p-value; оценивать новую политику на логах без propensity; игнорировать организационную стоимость MAB. Именно эти риски проверяются в экспериментальной части.

\textbf{Таблица 2.1 --- Сводка компромисса A/B и MAB}

\begingroup\footnotesize\setlength{\tabcolsep}{4pt}
\begin{xltabular}{\linewidth}{@{}
  >{\raggedright\arraybackslash}Y
  >{\raggedright\arraybackslash}Y
  >{\raggedright\arraybackslash}Y@{}}
\toprule\noalign{}
\midrule\noalign{}
\endhead
\bottomrule\noalign{}
\endlastfoot
Цель & Статистический вывод об эффекте & Максимизация награды при обучении \\
Стоимость внедрения & Ниже & Выше (организационно) \\
Интерпретируемость & Выше & Ниже \\
Офлайн-оценка & Прямое сравнение групп & OPE с propensity \\
Масштаб по числу рук & Хуже при росте K & Лучше при многоруких сценариях \\
\end{xltabular}
\endgroup

\subsection{Модель многорукого бандита и её основные идеи}\label{ux43cux43eux434ux435ux43bux44c-ux43cux43dux43eux433ux43eux440ux443ux43aux43eux433ux43e-ux431ux430ux43dux434ux438ux442ux430-ux438-ux435ux451-ux43eux441ux43dux43eux432ux43dux44bux435-ux438ux434ux435ux438}

Модель многорукого бандита пришла из области обучения с подкреплением. Её название связано с образом игрового автомата: несколько автоматов с рычагами дают выигрыш с неизвестной вероятностью, и игроку нужно последовательно выбирать рычаг, чтобы максимизировать суммарный выигрыш \cite{ref1}, \cite{ref2}. В прикладной постановке вместо автоматов --- варианты интерфейса, рекламные объявления, рекомендации товаров; вместо выигрыша --- клик, конверсия или другая наблюдаемая награда.

На каждом шаге алгоритм выбирает одно действие из конечного набора, получает награду, но заранее не знает лучшее действие. Возникает базовая дилемма exploration--exploitation: исследование менее изученных вариантов vs использование текущего лучшего. Слишком раннее exploitation закрепляет неоптимальный вариант; слишком долгое exploration накапливает потери от показа слабых альтернатив \cite{ref2}, \cite{ref14}, \cite{ref16}, \cite{ref17}.

В литературе выделяют три базовых подхода. \textbf{$\varepsilon$-greedy} с вероятностью $\varepsilon$ выбирает случайное действие, иначе --- вариант с наилучшей текущей оценкой \cite{ref13}, \cite{ref16}, \cite{ref17}. Метод прозрачен, но уровень исследования задаётся внешним параметром и не адаптируется к степени неопределённости. \textbf{UCB} (Upper Confidence Bound) добавляет малоизученным вариантам бонус за неопределённость \cite{ref4}, \cite{ref14}, \cite{ref17}: исследование определяется объёмом накопленной информации, а не фиксированной $\varepsilon$. \textbf{Thompson Sampling} поддерживает для каждого действия распределение неопределённости и на каждом шаге выбирает вариант, лучший в случайной выборке из этих распределений \cite{ref1}, \cite{ref8}, \cite{ref9}, \cite{ref16}. Исследование и использование возникают естественно из ширины распределения; на бинарной награде TS часто показывает сильные результаты \cite{ref8}, \cite{ref9}, \cite{ref16}.

Важнейшее понятие --- regret, накопленное отставание от стратегии, которая всегда выбирала бы наилучшее действие \cite{ref3}, \cite{ref14}, \cite{ref16}, \cite{ref17}. Для A/B-тестирования централен вопрос о наличии эффекта и его статистической значимости; для MAB --- минимизация потерь в процессе выбора. Модель изначально разрабатывалась как последовательное принятие решений в условиях неопределённости, а не как разновидность статистического эксперимента в привычном смысле; позднее получила широкое применение в онлайн-экспериментах, рекомендациях и рекламе \cite{ref2}, \cite{ref12}, \cite{ref13}, \cite{ref16}.

\subsection{Контекстные бандиты}\label{ux43aux43eux43dux442ux435ux43aux441ux442ux43dux44bux435-ux431ux430ux43dux434ux438ux442ux44b}

Следующий шаг после классического MAB --- контекстные бандиты (contextual bandits). Перед выбором действия алгоритм наблюдает признаки ситуации: характеристики пользователя, историю взаимодействий, тип устройства, позицию объекта. Награда наблюдается только по выбранному варианту \cite{ref6}, \cite{ref15}, \cite{ref16}, \cite{ref17}. В отличие от обычного бандита, где ищется один лучший вариант в среднем, контекстный бандит учится выбирать лучшее действие для конкретного наблюдаемого случая.

Контекстный бандит --- промежуточная модель между классическим MAB и полным RL: нет длинной цепочки состояний, но нужно изучить зависимость между признаками и ожидаемой наградой \cite{ref15}, \cite{ref16}, \cite{ref17}. Линейные методы (LinUCB) предполагают линейную зависимость награды от признаков \cite{ref6}, \cite{ref15}; более гибкие аппроксимации повышают сложность реализации \cite{ref16}, \cite{ref17}.

Практическая ценность особенно заметна в рекомендациях, рекламе и ранжировании \cite{ref6}, \cite{ref31}, \cite{ref37}. Каноническая работа Li et al.~\cite{ref6} показала контекстные бандиты как рабочий механизм персонализации в крупной системе и связала тему с логами и автономной оценкой на исторических данных. Офлайн A/B для рекомендательных систем обсуждается отдельно \cite{ref37}.

Для курсовой контекстные бандиты --- важное теоретическое расширение и естественное продолжение темы (RQ5, LinUCB), а основной объём базовой эмпирики строится на классических алгоритмах MAB и \texttt{fixed\_ab} \cite{ref13}, \cite{ref14}, \cite{ref16}, \cite{ref17}.

\subsection{Сравнение A/B-тестирования и MAB в современной литературе}\label{ux441ux440ux430ux432ux43dux435ux43dux438ux435-ab-ux442ux435ux441ux442ux438ux440ux43eux432ux430ux43dux438ux44f-ux438-mab-ux432-ux441ux43eux432ux440ux435ux43cux435ux43dux43dux43eux439-ux43bux438ux442ux435ux440ux430ux442ux443ux440ux435}

Работы прямого сравнения особенно важны для настоящей темы. В обзоре Burtini и соавторов \cite{ref13} MAB рассматривается как способ организации онлайн-эксперимента, где выигрыш от адаптивного распределения трафика сопоставляется со сложностью статистического вывода и требованиями внедрения. Scott \cite{ref12} связывает интерес к бандитам с экономикой цифровых сервисов: стоимость эксперимента определяется не только длительностью теста, но и потерями от показа плохого варианта реальным пользователям.

Литература подчёркивает и обратную сторону. Mattos, Bosch и Olsson \cite{ref18} показывают, что bandit-методы усложняют объяснение результатов, требуют зрелости экспериментальной платформы и не всегда вписываются в процессы, где ожидают привычный показатель эффекта и простое решение о победителе. Бандиты выигрывают прежде всего там, где цена ошибки высока, награда наблюдается быстро, горизонт достаточно длинный, а организация готова к адаптивному назначению трафика \cite{ref12}, \cite{ref13}, \cite{ref18}.

Xiang и соавторы \cite{ref21}, \cite{ref22} показывают, что адаптивные алгоритмы могут выигрывать у A/B по накопленной награде и regret, но выигрыш зависит от постановки и не универсален. Bleeker и Kaufmann \cite{ref22} рассматривают MAB как способ сделать эксперименты дешевле, а не как полную замену A/B. Корректная формулировка: MAB --- адаптивная альтернатива фиксированному распределению трафика там, где важен не только итоговый вывод, но и качество решений на протяжении обучения \cite{ref12}, \cite{ref13}, \cite{ref18}, \cite{ref21}, \cite{ref22}.

\subsection{Датасеты и эмпирические основания}\label{ux434ux430ux442ux430ux441ux435ux442ux44b-ux438-ux44dux43cux43fux438ux440ux438ux447ux435ux441ux43aux438ux435-ux43eux441ux43dux43eux432ux430ux43dux438ux44f}

Выбор данных требует отдельного обоснования: не всякий датасет рекомендаций позволяет корректно сравнивать фиксированное и адаптивное распределение трафика. Сильные примеры off-policy evaluation связаны с логами, где записаны не только действия и награды, но и propensity исходной политики \cite{ref6}, \cite{ref24}. Крупные recommendation-датасеты вроде MIND \cite{ref25} полезны для ранжирования и офлайн-метрик, но не подходят для IPS/SNIPS replay-OPE в постановке курсовой: в публичной версии нет полей propensity политики логирования, необходимых для IPS/SNIPS. Поэтому для внешней проверки выбран OBD \cite{ref26} с random и BTS logging. Дополнительно используется синтетическая Bernoulli-среда для контролируемого измерения regret.

Синтетика позволяет задать истинные CTR, точно вычислять regret, менять горизонт и расстояние между вариантами. OBD даёт реальное пользовательское поведение, propensity и многорукую постановку (80 товаров в кампании all). Двухуровневая эмпирика убедительнее, чем попытка сразу построить промышленную систему \cite{ref17}, \cite{ref21}, \cite{ref24}.

В курсовой дополнительно рассматривается проекция OBD на pairwise-постановку классического A/B (два варианта) (RQ7, E13): два популярных \texttt{item\_id} трактуются как условные варианты A и B. Это методологический мост между многоруким каталогом и продуктовым A/B, а не замена отдельного эксперимента с фиксированным разбиением трафика.

\subsection{Выводы из обзора и связь с экспериментальной частью}\label{ux432ux44bux432ux43eux434ux44b-ux438ux437-ux43eux431ux437ux43eux440ux430-ux438-ux441ux432ux44fux437ux44c-ux441-ux44dux43aux441ux43fux435ux440ux438ux43cux435ux43dux442ux430ux43bux44cux43dux43eux439-ux447ux430ux441ux442ux44cux44e}

Проведённый анализ позволяет сформулировать несколько выводов. В реальной эксплуатации уже существуют отработанные схемы онлайн-экспериментов: fixed A/B для статистически корректной проверки гипотезы \cite{ref5}, \cite{ref10}, \cite{ref40}, MAB и contextual bandits для адаптивного назначения в рекомендациях \cite{ref6}, \cite{ref12}, \cite{ref37}, OPE-пайплайны для оценки по логам \cite{ref24}, \cite{ref26}. Статистически значимый результат в индустрии получают не «из cumulative CTR бандита», а через фиксированный дизайн, планирование мощности, снижение дисперсии и --- при необходимости мониторинга --- последовательные процедуры \cite{ref7}, \cite{ref19}. Слабое место практики --- интерпретация логов: смешивание оптимизации и статистического вывода, наивные p-value на adaptive-назначении, подглядывание без поправки, OPE при низком ESS \cite{ref11}, \cite{ref18}, \cite{ref21}.

A/B-тестирование и MAB нельзя рассматривать как полностью взаимозаменяемые методы: A/B остаётся предпочтительным там, где нужен прозрачный и статистически корректный вывод об эффекте, MAB --- в задачах, где важно уменьшать потери уже в ходе эксперимента \cite{ref5}, \cite{ref12}, \cite{ref18}, \cite{ref19}, \cite{ref20}.

Для курсовой достаточно сильным ядром являются классические стохастические алгоритмы --- $\varepsilon$-greedy, UCB1 и Thompson Sampling --- в сопоставлении с \texttt{fixed\_ab} \cite{ref4}, \cite{ref8}, \cite{ref9}, \cite{ref13}, \cite{ref17}. Контекстные бандиты раскрыты в теоретической части и проверены в RQ5, но не являются центральным объектом базовой реализации \cite{ref6}, \cite{ref15}, \cite{ref17}.

Экспериментальная часть (глава 3) сфокусирована на онлайн-эксперименте в двух связанных блоках. Первый блок --- доказать на контролируемых и продуктовых постановках, что адаптивное назначение трафика (MAB, contextual bandit) приносит измеримую пользу по regret и накопленным кликам (E1, E6, E8b, E12). Второй блок --- показать, что стандартные способы получить «статистически значимый» вывод дают сбой на логах таких экспериментов, и проверить известные из литературы обходные пути: отдельный fixed A/B, IPS-weighted inference, group sequential и корректный мониторинг при непрерывном наблюдении (always-valid, см. сокращения) (E4, E5, E11, E14). Офлайн-оценка по логам OBD (E2, E9, E13) дополняет картину интерпретации уже собранных онлайн-экспериментов. Структура: главы 1--2 --- постановка, RQ (§1.6) и обзор; глава 3 --- эмпирика по четырём типам оценки (§1.8, §3.2), включая pairwise-постановку RQ7 (§3.6.4).
